\section{Selecting compatible translates}\label{sec:selection}

The deletion estimates require independent random choices, while
intersection of the exceptional edges requires geometric compatibility.
We separate these tasks by sampling several independent translates of
each bounded configuration and then selecting one from each pool.
The theorem below supplies the selection step. Its input consists of
finite point sets $W_i$ and sets $D_i$ of forbidden join directions;
its output avoids every forbidden join between different selected sets.
Because a direction is a linear subspace containing zero, the condition
also rules out a common point between two selected configurations
whenever at least one of their direction sets is nonempty.

\begin{theorem}[Compatible translates from independent pools]\label{thm:selection}
For every positive integer $D$ there is a positive integer $K=K(D)$
with the following property.  Let $q$ tend to infinity through prime
powers.  For each $q$, let $n\le Dq$, and for $1\le i\le n$ let
$W_i\subseteq\F_q^2$ and let $D_i$ be a set of directions, with
\[
 |W_i|\le D,\qquad |D_i|\le D.
\]
Assume that each direction belongs to $D_i$ for at most $D$ indices $i$.
For distinct indices $i,j$, declare the translates $b+W_i$ and
$b'+W_j$ to conflict if there are $u\in W_i$, $v\in W_j$, and
$d\in D_i\cup D_j$ such that
\[
 (b+u)-(b'+v)\in d.
\]
For each $i$, draw $K$ independent uniform shifts
$b_{i,1},\ldots,b_{i,K}\in\F_q^2$, independently for all indices.
With probability tending to one, one can choose $k_i\in\{1,\ldots,K\}$
for every $i$ so that the translates $b_{i,k_i}+W_i$ have no conflicts.
The convergence is uniform over all data satisfying the stated bounds.
\end{theorem}

To prove the theorem, make each candidate a vertex of a graph and
join candidates from different pools when they conflict. We seek an
independent set containing one vertex from each pool, called an
independent transversal. Its nonexistence has a small domination
certificate. We first prove the form of that assertion we need, then
show that the independent shifts make every such certificate unlikely.
The certificate argument is a weaker form of Haxell's
independent-transversal domination theorem~\cite{Haxell2001},
whose method originates in~\cite{Haxell1995}. We include its proof.
Haxell's theorem uses total domination, which also requires neighbors
for the vertices of the dominating set. Here we use the following
weaker convention.
For a graph $G$ and a set $S$ of its vertices, we say that $T\subseteq S$
\emph{dominates $S$} if every vertex of $S\setminus T$ has a neighbor in
$T$.  No condition is imposed on the vertices of $T$ themselves.

\begin{lemma}[A domination certificate]\label{sel-domination}
Let a finite simple graph have a vertex partition into nonempty parts
$A_1,\ldots,A_n$.  If there is no independent set containing one
vertex from each part, then there are a nonempty index set $I$ and a set
\[
 T\subseteq\bigcup_{i\in I}A_i,
 \qquad |T|\le 2(|I|-1),
\]
such that $T$ dominates $\bigcup_{i\in I}A_i$.
\end{lemma}

\begin{proof}
We grow a set of vertices with small size compared with the number
of parts it reaches. Whenever it does not yet dominate those parts,
we use an undominated vertex to reach another part. A minimality
condition on the transversals prevents this growth from stopping
without the required domination certificate.

Choose an inclusion-minimal nonempty subfamily of parts with no
independent transversal, and distinguish one part $A_*$ of this
subfamily.  Let $\mathcal U$ be the family of independent transversals of
all its other parts.  Minimality makes $\mathcal U$ nonempty.  The
deficient subfamily has at least two parts, since every part is nonempty.

Choose any $x_1\in A_*$.  Among $U\in\mathcal U$, minimize
$|N(x_1)\cap U|$, where $N(x)$ denotes the neighborhood of $x$ in $G$.
Fix a minimizer and write $Y_1=N(x_1)\cap U$.  This set is nonempty:
otherwise adjoining $x_1$ would give an independent transversal of the
deficient subfamily.

Inductively, suppose that vertices $x_1,\ldots,x_j$ and sets
$Y_1,\ldots,Y_j$ have been fixed with the following properties:
\begin{enumerate}
 \item The sets $Y_a$ are nonempty and pairwise disjoint, and some
 $U\in\mathcal U$ jointly realizes all the equalities
 \[
   N(x_a)\cap U=Y_a\qquad(1\le a\le j).
 \]
 \item At stage $a$, the size $|Y_a|$ was minimized over all
 $U\in\mathcal U$ satisfying the already fixed equalities for indices
 smaller than $a$.
 \item Each $x_a$ lies either in $A_*$ or in a part represented by a
 vertex of some earlier $Y_b$, with $b<a$.
\end{enumerate}
Let $Y=\bigcup_{a=1}^jY_a$, let $I$ consist of the distinguished part
and all parts represented in $Y$, and put
\[
 T=\{x_1,\ldots,x_j\}\cup Y.
\]
By joint realization, the vertices of $Y$ occupy distinct parts, none
of them $A_*$.  Thus $|I|=|Y|+1$, while nonemptiness of the $Y_a$ gives
$j\le |Y|$.  The third invariant shows that $T$ lies in the union of
the parts indexed by $I$, and
\[
 |T|\le j+|Y|\le 2|Y|=2(|I|-1).
\]
If $T$ dominates this union, we have the required certificate.

Otherwise, choose a vertex $x_{j+1}$ in that union, outside $T$, with
no neighbor in $T$.  Among the transversals realizing all the current
equalities, choose $U$ minimizing $|N(x_{j+1})\cap U|$, and set
\[
 Y_{j+1}=N(x_{j+1})\cap U.
\]
This set avoids all the previous $Y_a$, because $x_{j+1}$ has no
neighbor in $T$.  We claim it is nonempty.

If it were empty and $x_{j+1}\in A_*$, then $U\cup\{x_{j+1}\}$
would be an independent transversal of the deficient subfamily.
Otherwise, the part of $x_{j+1}$ contains a unique vertex $y$ of $Y$,
say $y\in Y_b$.  The transversal $U$ uses $y$ in that part, so
$x_{j+1}\notin U$.  Replace $y$ by $x_{j+1}$ to form
\[
 U'=(U\setminus\{y\})\cup\{x_{j+1}\}.
\]
Since $N(x_{j+1})\cap U=\varnothing$, this is again an independent
transversal of the nondistinguished parts.  For every $a<b$, the vertex
$y$ does not belong to $Y_a$, and the new vertex is not adjacent to
$x_a$.  Consequently the exact earlier equalities are preserved:
\[
 N(x_a)\cap U'=Y_a\qquad(a<b).
\]
For $a=b$, the replacement removes the neighbor $y$ and adds no neighbor,
so
\[
 |N(x_b)\cap U'|=|Y_b|-1.
\]
This contradicts the minimum imposed at stage $b$.  Notice that the
contradiction uses only the equalities preceding stage $b$; preservation
of later equalities is not required.

Thus $Y_{j+1}$ is nonempty and all the induction invariants continue to
hold.  Each continuation adds vertices in previously unrepresented
parts to $Y$, whereas all of $Y$ lies in one transversal of the fixed
nondistinguished subfamily.  There can be only finitely many
continuations.  The process must therefore stop with a dominating set
$T$, proving the lemma.
\end{proof}

\begin{proof}[Proof of Theorem~\ref{thm:selection}]
The case $n=0$ is immediate, so suppose $n\ge1$.  Fix $D$ and, for the
moment, a constant integer $K\ge6$.  We will choose $K$ sufficiently
large in terms of $D$ at the end.

Form a graph whose vertices are the candidate positions $(i,k)$,
partitioned into the $n$ pools.  Put an edge between candidates in
different pools precisely when their translates conflict.  Candidate
positions are distinct vertices even when their sampled shifts happen
to coincide.  An independent transversal in this graph is exactly a
selection as in the theorem.

By Lemma~\ref{sel-domination}, it suffices to show that, with probability
tending to one, there are no nonempty index sets $I$, with $|I|=m$, and
sets $T$ of at most $2m$ candidate positions in those pools that dominate
their union.  We use the slightly weaker bound $2m$ to simplify the
counting.  Throughout the probability estimates, $I$ and $T$ are fixed
sets of indices before any shifts are exposed.

The proof splits according to $m$, the number of pools in the
certificate. When $m$ is small, the forbidden lines associated with
$T$ occupy only a small part of each candidate's possible shifts.
When $m$ is proportional to $q$, that direct bound is insufficient.
Instead, we expose $T$ in a bounded number of chunks and prove that
a positive average proportion of shifts remains feasible. Crucially,
this proportion will depend only on $D$, so we may subsequently choose
$K$ large enough that the unused candidates defeat all certificates.

\paragraph{Forbidden equations and counting possible certificates.}
For every direction $d$ in use, fix a nonzero linear functional
$\pi_d:\F_q^2\to\F_q$ whose kernel is $d$.  After exposing the shifts
of the candidates $e\in T$, call a shift $b\in\F_q^2$ \emph{feasible
for target $i$} if, for every such $e$, with source-pool index $j$, it
avoids every equation
\begin{equation}\label{sel-forbidden}
 \pi_d(b)=\pi_d(b_e)+\pi_d(v-u),
 \qquad
 u\in W_i,\quad v\in W_j,\quad d\in D_i\cup D_j.
\end{equation}
We impose these prohibitions even when $i=j$.  This only makes
feasibility more restrictive.  A feasible candidate outside $T$ has no
neighbor in $T$, so its existence prevents domination.

Set
\begin{equation}\label{sel-constants}
 C=2D^3,\qquad c_1=\frac1{4C},\qquad \alpha=\frac1{2C}.
\end{equation}
For one target $i$ and one exposed candidate $e$, there are at most $C$
equations in \eqref{sel-forbidden}.  This remains true if the two lists
$D_i$ and $D_j$ are counted separately, retaining duplicate directions.
Each equation excludes one affine line, containing exactly $q$ shifts.

The number of choices of $I$ is at most
\[
 \binom{n}{m}\le\left(\frac{eDq}{m}\right)^m.
\]
There is an absolute constant $C_*>0$ such that, for $K\ge6$, the number
of candidate-index sets $T$ in these pools with $|T|\le2m$ is at most
\begin{equation}\label{sel-certificate-count}
 \sum_{t=0}^{2m}\binom{Km}{t}\le (C_*K^2)^m.
\end{equation}
For example, the usual bound
$\sum_{t=0}^{r}\binom Nt\le(eN/r)^r$ for $1\le r\le N/2$
applies with $N=Km$ and $r=2m$.

\paragraph{The range $1\le m\le c_1q$.}
Condition on all the exposed shifts in $T$.  A fresh uniform shift for
any target $i\in I$ is infeasible with probability at most
\[
 \frac{C|T|}{q}\le\frac{2Cm}{q}\le\frac12.
\]
At least $(K-2)m$ candidate positions in the target pools remain
unexposed.  Their shifts are conditionally independent and uniform.
Thus the probability that $T$ dominates, for this fixed $I,T$, is at most
\begin{equation}\label{sel-small-probability}
 \left(\frac{2Cm}{q}\right)^{(K-2)m}.
\end{equation}
Combining this with the counts of $I$ and $T$, the probability of a
dominating certificate of size parameter $m$ is at most
\begin{equation}\label{sel-small-sum}
 \left[
  2CeD C_*K^2
  \left(\frac{2Cm}{q}\right)^{K-3}
 \right]^m.
\end{equation}
Choose $K$ large enough that
\begin{equation}\label{sel-small-K}
 2CeD C_*K^2\,2^{-(K-3)}\le\frac12.
\end{equation}
Then the bracket in \eqref{sel-small-sum} is at most $1/2$ throughout
this range, and tends to zero for every fixed positive $m$.  Summing
over $m$ gives $o(1)$: for a fixed cutoff, the finite initial sum tends
to zero, while the remaining terms are bounded by a geometric tail.

\paragraph{A feasible-density estimate for $m>c_1q$.}
We prove that there is a constant $\delta_1>0$, depending only on $D$
and not on $K$, such that for every fixed $I,T$ in this range,
\begin{equation}\label{sel-density-target}
 \Pr\!\left[
   \sum_{i\in I}|U_i^{\mathrm{fin}}|<\delta_1mq^2
 \right]
 \le \exp\bigl(-\Omega_D(q\log q)\bigr).
\end{equation}
Here $U_i^{\mathrm{fin}}$ is the set of shifts feasible for target $i$
after all candidates in $T$ have been exposed.  The estimate is uniform
in $I,T$ and the configuration data.  The lower cutoff on $q$ may
depend on the fixed value of $K$.

Expose the candidate positions of $T$ in a fixed order, divided into
consecutive chunks of size at most $\lfloor\alpha q\rfloor$, with all
but the last having that size.  For sufficiently large $q$,
$\lfloor\alpha q\rfloor\ge\alpha q/2$.  Since $|T|\le2m\le2Dq$,
the number of chunks is at most
\begin{equation}\label{sel-chunk-count}
 J_0=\left\lceil\frac{4D}{\alpha}+1\right\rceil.
\end{equation}
In particular, this bound does not involve $K$.  If $T$ is empty, the
desired feasible-density estimate holds deterministically; otherwise
there is at least one chunk.

At the start of a chunk, condition on all shifts previously exposed.
Let $U_i$ be the set of shifts still feasible for target $i$, and write
\[
 M=\sum_{i\in I}|U_i|.
\]
For the analysis of this chunk we keep all these sets frozen at their
start values.  They are now deterministic, while the shifts belonging
to the new chunk are independent uniform vectors.

For an entry $e$ of the chunk, with source-pool index $j$, define $Z_e$
by summing, over all equations \eqref{sel-forbidden} associated with
$i\in I$ and $e$, the number of shifts $b\in U_i$ satisfying that
equation.  Count the $D_i$ and $D_j$ lists separately, including
duplicates.  A shift deleted more than once is therefore counted more
than once, which is harmless for an upper bound on loss.  Each line
contains at most $q$ shifts, and each of its right-hand sides is uniform
in $\F_q$.  Consequently,
\begin{equation}\label{sel-Z-basic}
 0\le Z_e\le Cmq,
 \qquad \E Z_e\le \frac{CM}{q}.
\end{equation}
All expectations and variances in the following chunk calculation are
conditional on the fixed past.

\paragraph{The terms whose directions belong to the target.}
The expectation bound alone does not control the loss in a chunk.
We estimate its fluctuations by separating the two direction lists
in each forbidden equation. Bounded reuse of target directions handles
one list; the affine second moment will handle the source list.
Write
\[
 Z_e=Z_e^{(1)}+Z_e^{(2)},
\]
where $Z_e^{(1)}$ uses the list $D_i$ in each target, and $Z_e^{(2)}$
uses the list $D_j$ of the source.  Group the terms of $Z_e^{(1)}$ by
direction.  The group for $d$ is a function only of $\pi_d(b_e)$.
For distinct directions $d,d'$, the functionals $\pi_d,\pi_{d'}$ are
linearly independent, so
\[
 b_e\longmapsto\bigl(\pi_d(b_e),\pi_{d'}(b_e)\bigr)
\]
is a bijection from $\F_q^2$ to itself.  The two coordinates are
therefore independent and uniform.  The direction groups are pairwise
independent, which suffices for adding their variances.

There are at most $Dm$ such groups.  Each direction belongs to at most
$D$ target direction sets; for each associated target-source pair
there are at most $D^2$ choices of $(u,v)$.  Thus every group is bounded
by $D^3q$, and
\begin{equation}\label{sel-first-variance}
 \operatorname{Var}(Z_e^{(1)})\le D^7mq^2.
\end{equation}

\paragraph{The terms whose directions belong to the source.}
For each target $i$ and each direction $d$, define the variance of its
parallel line slices by
\begin{equation}\label{sel-slice-variance}
 V_{i,d}=\frac1q\sum_{\ell\parallel d}
 \left(|U_i\cap\ell|-\frac{|U_i|}{q}\right)^2,
\end{equation}
where the sum is over the $q$ affine lines of direction $d$.
These variances obey the exact identity
\begin{equation}\label{sel-energy}
 \sum_d V_{i,d}
   =|U_i|-\frac{|U_i|^2}{q^2}
   \le |U_i|\le q^2,
\end{equation}
where the sum runs over all $q+1$ directions. This is precisely the
normalized affine second-moment identity
in~\cite[Lemma~A.1]{OpenAI2026Balanced}, applied to $U_i\subseteq\F_q^2$.
That identity holds for every subset of every finite affine plane;
it has no randomness, characteristic, or subfield hypothesis. Thus it
applies after conditioning on any history of earlier chunks. This is
the reason we can control all source directions even when the feasible
sets $U_i$ have complicated shapes.

Every individual summand of $Z_e^{(2)}$ associated with $i,d,u,v$ is a
uniformly sampled line-slice count, with its slice index translated by
a fixed offset.  Its variance is therefore $V_{i,d}$.  There are at
most $D^3m$ summands, and each pair $(i,d)$ is repeated at most $D^2$
times.  Cauchy--Schwarz applied to their centered sum, without any
independence assumption, gives
\begin{equation}\label{sel-second-variance}
 \operatorname{Var}(Z_e^{(2)})
 \le D^5m A_e,
 \qquad
 A_e:=\sum_{i\in I}\sum_{d\in D_j}V_{i,d}.
\end{equation}

Each source pool contributes at most $K$ entries to the chunk, and
each direction belongs to at most $D$ source direction sets.
Consequently, by \eqref{sel-energy},
\begin{equation}\label{sel-exception-energy}
 \sum_{e\text{ in the chunk}} A_e
 \le KD\sum_{i\in I}\sum_d V_{i,d}
 \le KDmq^2.
\end{equation}
Call an entry \emph{regular} if $A_e\le mq^{3/2}$.  There are at most
$KD\sqrt q$ irregular entries.  For a regular entry,
\eqref{sel-first-variance}, \eqref{sel-second-variance}, and
$\operatorname{Var}(X+Y)\le2\operatorname{Var}(X)+2\operatorname{Var}(Y)$
give
\begin{equation}\label{sel-regular-variance}
 \operatorname{Var}(Z_e)
 \le 2D^7mq^2+2D^5m^2q^{3/2}.
\end{equation}
For a fixed constant $\gamma>0$, Chebyshev's inequality and
\eqref{sel-Z-basic} therefore imply
\begin{align}
 \Pr\!\left[Z_e>\frac{CM}{q}+\gamma mq\right]
 &\le \frac{2D^7}{\gamma^2m}
       +\frac{2D^5}{\gamma^2\sqrt q}
 \notag\\
 &\le B_{D,\gamma}q^{-1/2}
 \qquad\text{for regular }e,
 \label{sel-regular-tail}
\end{align}
where $B_{D,\gamma}$ is independent of $K$.  In the last step we used
$m>c_1q$.

\paragraph{Conditional independence and loss over one chunk.}
The preceding variance bounds show that each regular entry has a
small chance of removing far more than its expected share. We now
use independence within the chunk to show that few entries do so,
and convert that statement into a deterministic bound on the total
remaining feasible set.

The regularity of an entry is determined by the frozen sets $U_i$ and
its fixed source-pool index.  In particular, it does not depend on its
fresh shift $b_e$.  After conditioning on the past, the collection of
regular entries is fixed, and their variables $Z_e$ are independent
functions of independent shifts.  Their common threshold in
\eqref{sel-regular-tail} is also fixed under this conditioning.

A union bound over subsets of regular entries, using the chunk size
bound and \eqref{sel-regular-tail}, shows that the probability of at
least $\gamma q$ threshold exceedances is at most
\begin{equation}\label{sel-chunk-tail}
 2^{\lceil\alpha q\rceil}
 \left(B_{D,\gamma}q^{-1/2}\right)^{\lceil\gamma q\rceil}
 \le \exp(-c_{D,\gamma}q\log q)
\end{equation}
for all sufficiently large $q$, with $c_{D,\gamma}>0$ independent of
$K$.  If there are fewer than $\lceil\gamma q\rceil$ regular entries,
the event in question is empty and the bound remains valid.

For each fixed $K$, we may also take $q$ sufficiently large that
\begin{equation}\label{sel-K-cutoff}
 KD\sqrt q\le\gamma q.
\end{equation}
Except on the event in \eqref{sel-chunk-tail}, all but at most
$2\gamma q$ entries are regular and below the threshold: the exceptions
include all irregular entries.  The actual loss of feasible shifts
within the chunk is at most $\sum_eZ_e$, because these counts used the
larger feasible sets at the start of the chunk.  Bounding each
exceptional entry by $Cmq$ and all other entries by the threshold gives,
for the end-of-chunk value $M'$, the deterministic inequality
\begin{align}
 M'
 &\ge M-C\alpha M-\alpha\gamma mq^2-2C\gamma mq^2
 \notag\\
 &=(1-C\alpha)M-(\alpha+2C)\gamma mq^2
 \notag\\
 &=\frac12M-(\alpha+2C)\gamma mq^2.
 \label{sel-chunk-recurrence}
\end{align}

Choose $\gamma>0$, depending only on $D$, so small that
\begin{equation}\label{sel-gamma-choice}
 2(\alpha+2C)\gamma\le 2^{-J_0-1},
 \qquad
 \delta_1:=2^{-J_0-1}.
\end{equation}
At the beginning, before exposing $T$, every shift is feasible, so
$M=mq^2$.  If all $J\le J_0$ chunks satisfy
\eqref{sel-chunk-recurrence}, iteration gives
\begin{align*}
 \frac{M_{\mathrm{fin}}}{mq^2}
 &\ge 2^{-J}-(\alpha+2C)\gamma
               \sum_{a=0}^{J-1}2^{-a}\\
 &\ge 2^{-J_0}-2(\alpha+2C)\gamma
 \ge\delta_1.
\end{align*}
The conditional probability bound \eqref{sel-chunk-tail} holds for
every possible past.  Averaging over the past and taking a union bound
over the at most $J_0$ chunks proves \eqref{sel-density-target}.
The constants $J_0$, $\gamma$, and $\delta_1$ depend only on $D$.
The dependence on $K$ in \eqref{sel-K-cutoff} affects only how large
$q$ must be, not these constants.

\paragraph{Unused candidates and the large-range union bound.}
We have established the required positive average feasible density
with a failure probability of order $\exp(-\Omega_D(q\log q))$.
There are only exponentially many candidate certificates, so this
error can be summed over all of them. On the successful event, the
unused candidates now supply the further exponential saving needed
to rule out domination.

Condition now on all the shifts belonging to $T$, and suppose the
feasible-density event in \eqref{sel-density-target} succeeds.  Let
$t_i$ be the number of exposed candidate positions from pool $i$.
Every unused candidate from that pool remains independent and uniform,
and is feasible with probability $|U_i^{\mathrm{fin}}|/q^2$.
The probability that none of the unused candidates in the target pools
is feasible is therefore at most
\begin{align}
 \prod_{i\in I}
 \left(1-\frac{|U_i^{\mathrm{fin}}|}{q^2}\right)^{K-t_i}
 &\le
 \exp\!\left(-\sum_{i\in I}
       (K-t_i)\frac{|U_i^{\mathrm{fin}}|}{q^2}\right)
 \notag\\
 &\le \exp\bigl(-(K\delta_1-2)m\bigr).
 \label{sel-unused-candidates}
\end{align}
To obtain the last inequality, use
\[
 \sum_{i\in I}\frac{|U_i^{\mathrm{fin}}|}{q^2}\ge\delta_1m,
 \qquad
 0\le\frac{|U_i^{\mathrm{fin}}|}{q^2}\le1,
 \qquad
 \sum_{i\in I}t_i=|T|\le2m.
\]
This argument also covers pools for which every candidate belongs to
$T$; it does not assume $K-t_i$ is positive in every target pool.

For $m>c_1q$, the number of pairs $I,T$ is at most $A_K^m$, where
\begin{equation}\label{sel-large-count}
 A_K=\frac{eD}{c_1}C_*K^2.
\end{equation}
For each fixed pair, domination can occur only if the feasible-density
estimate fails or all unused candidates are infeasible despite its
success.  The sum of the first failure probabilities over all
$m,I,T$ in the large range is bounded by
\[
 \exp\bigl(O_{D,K}(q)\bigr)
 \exp\bigl(-\Omega_D(q\log q)\bigr)=o(1).
\]
Here the number of choices is at most
$\sum_{m\le Dq}A_K^m=\exp(O_{D,K}(q))$, with $K$ fixed.

Finally choose $K$ sufficiently large, still depending only on $D$, so
that both \eqref{sel-small-K} and
\begin{equation}\label{sel-large-K}
 A_K\exp(2-K\delta_1)<\frac12
\end{equation}
hold.  Such a choice exists because $\delta_1>0$ does not depend on $K$,
whereas $A_K$ grows only quadratically in $K$.  By
\eqref{sel-unused-candidates}, the sum of the remaining failure
probabilities is then at most
\[
 \sum_{m>c_1q}
 \left(A_K\exp(2-K\delta_1)\right)^m
 \le\sum_{m>c_1q}2^{-m}=o(1).
\]
Together with the small-range estimate, this excludes all dominating
certificates with probability tending to one.  Lemma~\ref{sel-domination}
therefore supplies an independent transversal, proving the theorem.

The choices have been made in the order
\[
 D\ ;\quad C,c_1,\alpha,J_0,\gamma,\delta_1\ ;\quad K\ ;\quad q.
\]
In particular, $K$ is fixed before taking $q$ to infinity, and all
sufficiently-large-$q$ conditions may depend on that fixed $K$.
Every estimate above is uniform over the configurations and direction
sets subject to the bounds in the theorem.
\end{proof}
